% p.23 — الباب الثاني العشر في باع يبيع (تابع): المنفيات والمؤكدات
\mspage{23}
% Layout: six tall columns (right-to-left), each a stack of boxed blocks of unequal height.
% Every sub-heading box holds the heading AND, beneath it, the underlying (pre-iʿlāl) form,
% e.g. «مجهول بلم / لَمْ يُبْيَعْ»; the regular forms follow outside the box. Row alignment is
% approximate: all six columns run to 21 lines, so rows are aligned 1:1; the \cline marks
% reproduce the printed box boundaries inside each column.
% Column 1 continues the مضارع بلم (معروف) block begun on p.22.
\begin{sarf}{|c|c|c|c|c|c|}
لَمْ تَبِيْعَا & لَمْ أُبَعْ & لَنْ يُبَاعَا & لَتَبِيْعَنَّ & لَتُبَعْنَانِّ & \thead{أمر حاضر بالنون}\\
لَمْ تَبِيْعُوْا & لَمْ نُبَعْ & لَنْ يُبَاعُوْا & لَتَبِيْعَانِّ & لَأُبَاعَنَّ & اِبْيِعَنَّ\\\cline{2-2}\cline{6-6}
لَمْ تَبِيْعِي & \thead{مضارع بلن} & لَنْ تُبَاعَ & لَتَبِيْعُنَّ & لَنُبَاعَنَّ & بِيْعَنَّ\\\cline{5-5}
لَمْ تَبِيْعَا & لَنْ يَبْيِعَ & لَنْ تُبَاعَا & لَتَبِيْعِنَّ & \thead{منفي بالنون} & بِيْعَانِّ\\\cline{2-2}
لَمْ تَبِعْنَ & لَنْ يَبِيْعَ & لَنْ يُبَعْنَ & لَتَبِيْعَانِّ & لَا يَبْيِعَنَّ & بِيْعُنَّ\\\cline{5-5}
لَمْ أَبِعْ & لَنْ يَبِيْعَا & لَنْ تُبَاعَ & لَتَبِعْنَانِّ & لَا يَبِيْعَنَّ & بِيْعِنَّ\\\cline{5-5}
لَمْ نَبِعْ & لَنْ يَبِيْعُوْا & لَنْ تُبَاعَا & \com{لَا بِيْعَنَّ}{كذا في الأصل، بألف بلا همزة؛ ولعل الصواب: لَأَبِيْعَنَّ} & \thead{مجهول بالنون} & بِيْعَانِّ\\\cline{1-1}
\thead{مجهول بلم} & لَنْ تَبِيْعَ & لَنْ تُبَاعُوْا & لَنَبِيْعَنَّ & لَا يُبْيَعَنَّ & بِعْنَانِّ\\\cline{4-4}\cline{5-5}\cline{6-6}
لَمْ يُبْيَعْ & لَنْ تَبِيْعَا & لَنْ تُبَاعِي & \thead{مجهول بالنون} & لَا يُبَاعَنَّ & \thead{نهي حاضر بالنون}\\\cline{1-1}\cline{5-5}
لَمْ يُبَعْ & لَنْ يَبِعْنَ & لَنْ تُبَاعَا & لَيُبْيَعَنَّ & \thead{أمر غائب بالنون} & لَا تَبْيِعَنَّ\\\cline{4-4}\cline{6-6}
لَمْ يُبَاعَا & لَنْ تَبِيْعَ & لَنْ تُبَعْنَ & لَيُبَاعَنَّ & لِيَبْيِعَنَّ & لَا تَبِيْعَنَّ\\\cline{5-5}
لَمْ يُبَاعُوْا & لَنْ تَبِيْعَا & لَنْ أُبَاعَ & لَيُبَاعَانِّ & لِيَبِيْعَنَّ & لَا تَبِيْعَانِّ\\\cline{5-5}
لَمْ تُبَعْ & لَنْ تَبِيْعُوْا & لَنْ نُبَاعَ & لَيُبَاعُنَّ & \thead{مجهول بالنون} & لَا تَبِيْعُنَّ\\\cline{3-3}
لَمْ تُبَاعَا & لَنْ تَبِيْعِي & \thead{مضارع بالنون} & لَتُبَاعَنَّ & لِيُبْيَعَنَّ & لَا تَبِيْعِنَّ\\\cline{5-5}
لَمْ يُبَعْنَ & لَنْ تَبِيْعَا & لَيَبْيِعَنَّ & لَتُبَاعَانِّ & لِيُبَاعَنَّ & لَا تَبِيْعَانِّ\\\cline{3-3}\cline{5-5}
لَمْ تُبَعْ & لَنْ تَبِعْنَ & لَيَبِيْعَنَّ & لَيُبَعْنَانِّ & \thead{نهي غائب بالنون} & لَا تَبِعْنَانِّ\\\cline{6-6}
لَمْ تُبَاعَا & لَنْ أَبِيْعَ & لَيَبِيْعَانِّ & لَتُبَاعَنَّ & لَا يَبْيِعَنَّ & \thead{تمرين}\\\cline{5-5}\cline{6-6}
لَمْ تُبَاعُوْا & لَنْ نَبِيْعَ & لَيَبِيْعُنَّ & لَتُبَاعَانِّ & لَا يَبِيْعَنَّ & رَاعَ\\\cline{2-2}\cline{5-5}
لَمْ تُبَاعِي & \thead{مجهول بلن} & لَتَبِيْعَنَّ & لَتُبَاعُنَّ & \thead{مجهول بالنون} & شَاعَ\\
لَمْ تُبَاعَا & لَنْ يُبْيَعَ & لَتَبِيْعَانِّ & لَتُبَاعِنَّ & لَا يُبْيَعَنَّ & دَاخَ\\\cline{2-2}\cline{5-5}
لَمْ تُبَعْنَ & لَنْ يُبَاعَ & لَيَبِعْنَانِّ & لَتُبَاعَانِّ & لَا يُبَاعَنَّ & لَاعَ\\
\end{sarf}
